\documentclass[11pt]{article}
\usepackage[utf8]{inputenc}
\usepackage[T1]{fontenc}
\usepackage[margin=0.85in]{geometry}
\usepackage{newtxtext,newtxmath}
\usepackage{amsmath,array,calc}
\usepackage{graphicx}
\usepackage{booktabs}
\usepackage{longtable}
\usepackage{etoolbox}
\usepackage{hyperref}
\usepackage{enumitem}
\usepackage{caption}
\newcommand{\real}[1]{#1}
\providecommand{\tightlist}{\setlength{\itemsep}{0pt}\setlength{\parskip}{0pt}}
\AtBeginEnvironment{longtable}{\scriptsize}
\setlist{nosep}
\setlength{\parskip}{4pt}
\setlength{\parindent}{0pt}
\title{V-CARA: Uncertainty-Gated Counterfactual Reinforcement Learning for Visual Control on an Infinite Track}
\author{Nishant Sirohi}
\date{October 2026}
\begin{document}
\maketitle


\begin{abstract}

Visual control on an effectively unbounded track requires sustained
progress while balancing speed, stability, and off-track risk. We
present V-CARA, a visual reinforcement-learning system built on Soft
Actor-Critic (SAC) and augmented with an ensemble of latent transition
models, uncertainty-gated local counterfactual action evaluation, a risk
estimator, future-coverage prediction, and coverage-aware replay. The
agent consumes four stacked 64 by 64 grayscale frames and uses a
time-aware state representation. We analyze the seed-42
evaluation trajectory through 1,000,000 environment frames and a V-CARA
ablation with counterfactual evaluation disabled through 450,002 frames.
At one million frames, V-CARA records mean return 875.7 \(\pm\) 30.0,
mean coverage 0.985, median coverage 0.980, P90 coverage 0.998, off-road
fraction 0.090, mean speed 65.0, and T90 of 650 frames. In the common
training window, the no-CF variant learns high coverage earlier at some
checkpoints, while full V-CARA has substantially lower off-road fraction
and higher coverage around 250k--325k frames; the difference is not
uniform across all checkpoints. The results motivate reporting progress,
coverage distributions, off-road behaviour, speed, and time-to-threshold
jointly rather than relying on return alone.

\textbf{Keywords:} visual reinforcement learning; autonomous control;
Soft Actor-Critic; counterfactual action evaluation; latent dynamics;
uncertainty gating; risk-sensitive control.

\end{abstract}
\section{Introduction}\label{introduction}

Visual continuous control is difficult because observations are images
rather than direct access to the full physical state. The controller
must infer motion, track geometry, and the consequences of steering and
throttle from temporally correlated visual input. On an effectively
unbounded track, the objective is to sustain forward progress while
managing speed and avoiding off-road states.

A single aggregate reward can conceal important behaviour. A high-return
policy may be fast but unstable, while a conservative policy may reduce
off-road events at the expense of progress. Coverage is therefore
reported alongside return, median and upper-quantile coverage, off-road
fraction, speed, and time to 90\% coverage (T90). This paper describes
V-CARA and compares the full system with an ablation that disables local
counterfactual evaluation (no-CF).

The contribution is an integrated description of the algorithm, the
implementation corrections that motivated it, and a transparent analysis
of the recorded training trajectories. The paper does not interpret one
training seed as evidence of universal superiority. Instead, it
identifies where the two recorded runs differ, where they do not, and
which conclusions remain limited by the available evaluation design.

\section{Background and Related Work}\label{background-and-related-work}

V-CARA uses Soft Actor-Critic as its off-policy continuous-control
backbone. SAC optimizes expected return together with an entropy term,
encouraging a stochastic policy while learning from replay. The
maximum-entropy formulation is useful when the action consequences are
uncertain and exploration is important (Haarnoja et al., 2018).

Model-based reinforcement learning learns a predictive model and uses it
to evaluate candidate actions or generate learning signals. Such methods
can improve data efficiency, but their benefit depends on model fidelity
and on how prediction error compounds over rollout horizon (Janner et
al., 2019). V-CARA uses a short-horizon latent transition ensemble
rather than relying on long open-loop image prediction.

Ensembles provide a practical proxy for epistemic uncertainty:
disagreement among independently trained predictors can indicate states
or actions where the model is less reliable. This signal is not a
calibrated probability of failure; it is used as a gate that limits
counterfactual influence when disagreement is high. Replay-based
learning can also be sensitive to which experiences are sampled.
Prioritized replay established the value of non-uniform sampling based
on learning relevance (Schaul et al., 2016); V-CARA instead uses
progress-quality partitions intended to retain useful high-coverage and
elite transitions.

Deep reinforcement learning results can vary substantially across random
seeds and implementation details, making complete reporting important
(Henderson et al., 2018). Accordingly, the numerical findings below are
framed as observations from the recorded seed-42 trajectories, not as a
multi-seed estimate.

\section{Problem Formulation}\label{problem-formulation}

Let the environment be a partially observed control process. At time t,
the latent physical state is s\_t, the image observation is o\_t, and
the continuous action is a\_t. The policy receives a stack of recent
images and a time feature, then emits steering, throttle, and brake
controls. The environment advances with action repeat r=4, so a single
policy action is held for several simulator frames.

The objective is sustained progress under a finite episode horizon, with
penalties or termination rules for off-road behaviour, out-of-field
states, and stalls. Because the track is effectively infinite,
evaluation focuses on continuous progress rather than terminal-event
metrics. Coverage \(c_t\in[0,1]\) denotes normalized progress within the
evaluation horizon; it is an operational metric rather than a terminal
event. T90 is the number of episode frames needed to reach coverage
0.90; a value of -1 indicates that the threshold was not reached in the
recorded summary.

The primary evaluation vector is m=(R, C\_mean, C\_median, C\_P90, O, V,
T90), where R is mean return, C values summarize episode-level coverage,
O is off-road fraction, and V is mean speed. Lower O and lower T90 are
generally desirable; speed is interpreted jointly with coverage and
off-road behaviour rather than maximized without constraint.

\section{V-CARA Method}\label{v-cara-method}

\subsection{Overview}\label{overview}

V-CARA combines a visual SAC agent with five supporting mechanisms: (i)
a convolutional encoder over stacked grayscale images; (ii) a compact
latent-dynamics ensemble; (iii) local counterfactual action evaluation
with uncertainty gating; (iv) auxiliary risk and future-coverage
prediction; and (v) replay sampling that emphasizes progress and
high-quality trajectories. Time-aware inputs and explicit stall/timeout
handling address temporal and termination semantics.

\subsection{Visual encoder and
actor--critic}\label{visual-encoder-and-actorcritic}

The observation is resized to 64 by 64 grayscale and stacked over four
frames. A convolutional encoder maps the stack to a 128-dimensional
feature vector. An actor maps features to a distribution over continuous
actions. Twin critics estimate soft action values, and target critics
are updated by exponential moving average. The saved configuration uses
discount \(\gamma=0.99\), target update \(\tau=0.005\), actor learning
rate \(10^{-4}\), critic learning rate 2\(\times\)\(10^{-4}\), batch size 256,
replay capacity 300,000, and one gradient update per environment step.

\subsection{Latent transition
ensemble}\label{latent-transition-ensemble}

An ensemble of three predictors estimates the next latent feature given
the current feature and action. Let f\_t be the encoded state and g\_j
the j-th predictor. Each model predicts the next latent state from the
current feature and action. The ensemble mean is computed by averaging
the predicted next-state features: \[
\hat{\mu}_{t+1}=\frac{1}{K}\sum_{j=1}^{K}\hat{f}_{t+1}^{(j)}.
\] Disagreement is represented by a scalar u\_t, computed from
dispersion across predictions after normalization to a reference scale.
The uncertainty value is a model diagnostic and should not be
interpreted as a calibrated confidence probability.

\subsection{Counterfactual action
evaluation}\label{counterfactual-action-evaluation}

Given the actor action, V-CARA constructs a small local set of candidate
actions that includes the actor action and bounded perturbations around
it. Candidates are clipped to the action bounds and scored using
predicted value, risk, progress, and model disagreement. A generic
candidate score is: \[
S(a)=\bar{Q}(z,a)-\beta_r\hat{R}(z,a)+\lambda_p\hat{P}(z,a)-\lambda_u\hat{U}(z,a),
\] where Q-bar is the twin-critic aggregate, R-hat is predicted risk,
P-hat is predicted progress, and U-hat is model disagreement. The materials available for this manuscript do not specify numerical values
for $\beta_r$, $\lambda_p$, or $\lambda_u$. These values cannot be
recovered from aggregate evaluation summaries and are marked as not reported
rather than inferred. The evaluation CSV is a derived checkpoint summary, not
a substitute for the raw per-episode evaluation records.

Counterfactual influence is gated by model reliability. The gate combines
an uncertainty-threshold test, a warm-up eligibility flag, and an additional
quality gate $q_t$. The candidate-action rule below makes the decision
explicit: use the actor action if uncertainty exceeds the limit, eligibility
conditions fail, or the candidate does not exceed the actor score by the
required margin; otherwise use the selected candidate. The exact interaction
between the quality gate and action selection should be verified against the
experiment code before publication. In the saved
configuration, the uncertainty limit is 1.5, the counterfactual margin
is 0.02, target acceptance rate is 0.25, fixed horizon is 3, and
warm-up/ramp is 10,000 frames. A key V2 correction was to put
uncertainty and acceptance statistics on compatible scales so the
mechanism does not remain effectively inactive.

\subsection{Risk and progress
prediction}\label{risk-and-progress-prediction}

A risk head predicts near-term off-road or unsafe-state likelihood. A
progress head estimates future progress from latent features. These
auxiliary signals help distinguish candidate actions that obtain
short-term reward from those that preserve stable forward movement. The
risk objective is trained from available transition labels, while
progress prediction uses transition-local targets rather than a single
episode-wide label.

\subsection{Future-coverage
prediction}\label{future-coverage-prediction}

A future-coverage classifier estimates whether a state is likely to
reach a high-coverage threshold within a finite prediction window. The
label is attached to each transition according to its own future
trajectory, rather than propagating an episode-level success label to
all early states. A quality gate based on positive-label availability
and balanced accuracy prevents noisy early predictions from dominating
actor updates. This classifier is an auxiliary progress signal, not a
terminal-success detector.

\subsection{Coverage-aware replay}\label{coverage-aware-replay}

Replay is partitioned into ordinary, near-high-coverage, and elite
transitions. The saved mixture is 0.625 ordinary, 0.250
near-high-coverage, and 0.125 elite, subject to buffer availability.
Near-high-coverage begins at 0.75 coverage; elite coverage begins at
0.90 with maximum off-road fraction 0.15. The replay rule uses
trajectory quality and coverage, not only the environment's termination
flag. This corrects an earlier failure where useful time-limit or stall
trajectories were omitted from high-quality buffers.

\subsection{Action mapping and episode
semantics}\label{action-mapping-and-episode-semantics}

The wrapper uses action repeat 4, a 1,000-step time limit, a 160-frame
stall threshold, and a brake dead zone of 0.6. Weak brake outputs are
mapped to zero to avoid unintended deceleration. The system
distinguishes genuine terminal events from time-limit truncations when
constructing critic targets. A stall mode applies an explicit penalty
schedule over up to 400 frames; no terminal progress bonus is used
because the track has no endpoint.

\subsubsection{Objective and uncertainty
gate}\label{objective-and-uncertainty-gate}

The SAC critic objective is \[
J_Q(\phi)=\mathbb{E}_{(s,a,r,s')\sim\mathcal D}
\left[\frac12\left(Q_\phi(s,a)-y\right)^2\right],
\] with soft target \[
y=r+\gamma(1-d)\left(\min_i Q_{\bar\phi_i}(s',a')-\alpha\log\pi(a'|s')\right).
\] For local counterfactual candidates \(a^{(j)}\), the score is \[
S(a^{(j)})=\bar Q(z,a^{(j)})-\beta_r\hat R(z,a^{(j)})
+\lambda_p\hat P(z,a^{(j)})-\lambda_u\hat U(z,a^{(j)}).
\] The counterfactual action is accepted only when the uncertainty and
eligibility gates pass and its score is at least the actor score plus the
margin $\delta$. Equivalently, the manuscript-level action rule is
\[
a_t=\begin{cases}
 a_{\pi}, & \hat U(z,a^{(j^*)})>u_{\max}\ \text{or eligibility fails, or}\ S(a^{(j^*)})<S(a_\pi)+\delta,\\
 a^{(j^*)}, & \text{otherwise}.
\end{cases}
\]
This rule is intended to match the prose in Section 4.4, but it has not been
verified against source code; the exact interaction with the quality gate
requires implementation-level confirmation.
The available materials do not state the perturbation magnitude or exact
candidate offsets, the risk-head calibration procedure, or all auxiliary-loss
coefficients. These are material reproducibility gaps: the aggregate results
cannot identify them. The reported uncertainty limit (1.5), counterfactual
margin (0.02), target acceptance rate (0.25), fixed horizon (3), and 10,000-frame
warm-up/ramp are documented, but they do not substitute for the missing values.

\subsection{Parameter disclosure and reproducibility status}\label{parameter-disclosure}

Table~\ref{tab:parameter-disclosure} distinguishes reported configuration values from parameters not specified in the method description. Missing values are not guessed. A reproducible release must include the exact run configuration or code commit that generated these logs.

\begin{longtable}{p{0.32\linewidth}p{0.19\linewidth}p{0.39\linewidth}}
\caption{Reported and missing implementation parameters.}\label{tab:parameter-disclosure}\\
\toprule
Parameter & Status & Value or note \\
\midrule
\endfirsthead
\caption[]{Reported and missing implementation parameters (continued).}\\
\toprule
Parameter & Status & Value or note \\
\midrule
\endhead
Observation size / stack & Reported & $64\times64$ grayscale; 4 frames \\
Latent ensemble size & Reported & 3 predictors \\
Discount / target update & Reported & $\gamma=0.99$, $\tau=0.005$ \\
Actor / critic learning rate & Reported & $10^{-4}$ / $2\times10^{-4}$ \\
Batch / replay capacity & Reported & 256 / 300,000 \\
Uncertainty limit / score margin & Reported & 1.5 / 0.02 \\
Target acceptance rate / horizon & Reported & 0.25 / 3 \\
Warm-up/ramp & Reported & 10,000 frames \\
Counterfactual perturbation magnitude and offsets & Not reported & Exact candidate action deltas unavailable \\
Score weights $\beta_r,\lambda_p,\lambda_u$ & Not reported & Numerical values unavailable \\
Risk calibration & Not reported & Calibration method, data split, and metric unavailable \\
Auxiliary-loss coefficients & Not reported & Complete coefficient list unavailable \\
Counterfactual acceptance counts by checkpoint & Not reported & Per-checkpoint gate/acceptance counts are not reported \\
\bottomrule
\end{longtable}

\section{Experimental Setup}\label{experimental-setup}

The full V-CARA evaluation series contains checkpoints from 100,000
through 1,000,000 environment frames. The separate no-CF seed-42 log
contains checkpoints from 25,000 through 450,002 frames. The no-CF
variant disables local counterfactual evaluation while retaining the
broader learning system. The runs share the reported seed identifier,
but because only one seed per variant is represented, differences are
descriptive rather than a reliable estimate of across-seed treatment
effects.

The full V-CARA table reports mean return with standard deviation, mean
coverage, median coverage, P90 coverage, off-road fraction, mean speed,
and T90. The no-CF log additionally reports exponential moving average
(EMA) return and out-of-field (OOF) fraction. The periodic
evaluation summaries are treated as the recorded values; no missing
checkpoints are interpolated for numerical tables.

For matched-window comparison, checkpoints are paired approximately by
training frames, with differences of a few frames preserved in the
source logs. This is not a paired statistical test: the metric is
evaluated from a single seed trajectory for each condition. A robust
confirmatory study would use multiple independent training seeds and a
fixed, shared evaluation seed set.

\section{Results}\label{results}

\subsection{Full V-CARA through one million
frames}\label{full-v-cara-through-one-million-frames}

At 1,000,000 frames, V-CARA records return 875.7 \(\pm\) 30.0, mean
coverage 0.985, median coverage 0.980, P90 coverage 0.998, off-road
fraction 0.090, mean speed 65.0, and T90=650. Mean coverage increases
from 0.527 at 100k to 0.985 at 1M, while off-road fraction falls from
0.110 at 100k to 0.090 at 1M. The trajectory is not monotonic: return and coverage deteriorate around
425k--500k, followed by partial recovery. At 700k, the change from 675k
is comparatively small: return is 866.0 \(\pm\) 45.0 versus 872.3
\(\pm\) 40.6, and mean coverage is 0.962 versus 0.965. The later
checkpoints recover and improve on these values.

From 750k onward, mean return fluctuates between 870.4 and 879.4 while
coverage rises from 0.975 to 0.985 and off-road fraction declines from
0.102 to 0.090. Across the late-training interval from 750k to 1M frames, mean return ranges from 870.4 to 879.4, while mean coverage increases from 0.975 to 0.985 and off-road fraction decreases from 0.102 to 0.090. Return remains within a relatively narrow range as coverage and off-road behaviour continue to change; these observations alone do not establish saturation of the return objective.

\subsection{V-CARA without counterfactual
evaluation}\label{v-cara-without-counterfactual-evaluation}

The no-CF trajectory reaches mean coverage 0.93 by 125k and records
return 838.4 \(\pm\) 58.1, median coverage 0.93, P90 0.97, off-road
fraction 0.05, and mean speed 52.2. This indicates faster early progress
than full V-CARA at the corresponding early checkpoints. The no-CF run
is also variable: at 250k, mean coverage falls to 0.80 with return 704.6
\(\pm\) 187.3 and off-road fraction 0.22; at 300k, mean coverage falls
to 0.75, return to 649.0 \(\pm\) 219.4, and off-road fraction rises to
0.28. By 450,002 frames it recovers to mean coverage 0.94, return 848.6
\(\pm\) 55.7, and off-road fraction 0.15.

\subsection{Common-window comparison}\label{common-window-comparison}

The comparison is phase-dependent. At approximately 100k--225k frames,
no-CF has higher mean coverage and return than full V-CARA in the
supplied runs. At 250k--325k, full V-CARA is more robust in several
checkpoints: at 300k, it records mean coverage 0.947 and off-road
fraction 0.106, versus 0.75 and 0.28 for no-CF; at 325k, it records
coverage 0.947 and off-road fraction 0.108, versus 0.79 and 0.24. At
350k--450k, coverage is similar, with no-CF marginally higher in some
checkpoints and full V-CARA has a lower off-road fraction at each reported checkpoint from 350k through 450k.

These observations do not support the blanket statement that
counterfactual evaluation improves every metric at every checkpoint.
They support a more specific interpretation: the counterfactual-equipped
run has a strong advantage during a mid-training instability window in
this seed-42 trajectory, while no-CF learns high coverage faster early
and later recovers. The one-million-frame full V-CARA endpoint is not
directly comparable to no-CF because the no-CF trajectory ends at
450,002 frames.

\begin{longtable}[]{@{}
  >{\raggedleft\arraybackslash}p{(\columnwidth - 14\tabcolsep) * \real{0.1250}}
  >{\raggedleft\arraybackslash}p{(\columnwidth - 14\tabcolsep) * \real{0.1250}}
  >{\raggedleft\arraybackslash}p{(\columnwidth - 14\tabcolsep) * \real{0.1250}}
  >{\raggedleft\arraybackslash}p{(\columnwidth - 14\tabcolsep) * \real{0.1250}}
  >{\raggedleft\arraybackslash}p{(\columnwidth - 14\tabcolsep) * \real{0.1250}}
  >{\raggedleft\arraybackslash}p{(\columnwidth - 14\tabcolsep) * \real{0.1250}}
  >{\raggedleft\arraybackslash}p{(\columnwidth - 14\tabcolsep) * \real{0.1250}}
  >{\raggedleft\arraybackslash}p{(\columnwidth - 14\tabcolsep) * \real{0.1250}}@{}}
\caption{Full V-CARA evaluation trajectory}\label{tab:full-vcara}\\
\toprule\noalign{}
\begin{minipage}[b]{\linewidth}\raggedleft
Frames
\end{minipage} & \begin{minipage}[b]{\linewidth}\raggedleft
Return \(\pm\) SD
\end{minipage} & \begin{minipage}[b]{\linewidth}\raggedleft
Mean cov.
\end{minipage} & \begin{minipage}[b]{\linewidth}\raggedleft
Median
\end{minipage} & \begin{minipage}[b]{\linewidth}\raggedleft
P90
\end{minipage} & \begin{minipage}[b]{\linewidth}\raggedleft
Off-road
\end{minipage} & \begin{minipage}[b]{\linewidth}\raggedleft
Speed
\end{minipage} & \begin{minipage}[b]{\linewidth}\raggedleft
T90
\end{minipage} \\
\midrule\noalign{}
\endhead
\bottomrule\noalign{}
\endlastfoot
100,000 & 427.0 \(\pm 134.7\) & 0.527 & 0.542 & 0.637 & 0.110 & 28.7 &
-1 \\
125,000 & 504.7 \(\pm 141.1\) & 0.605 & 0.632 & 0.717 & 0.123 & 33.7 &
217 \\
150,000 & 582.4 \(\pm 147.5\) & 0.683 & 0.723 & 0.797 & 0.137 & 38.8 &
435 \\
175,000 & 660.1 \(\pm 153.8\) & 0.760 & 0.813 & 0.877 & 0.150 & 43.8 &
653 \\
200,000 & 737.8 \(\pm 160.2\) & 0.838 & 0.903 & 0.957 & 0.163 & 48.8 &
871 \\
225,000 & 765.2 \(\pm 125.8\) & 0.865 & 0.913 & 0.962 & 0.149 & 50.2 &
865 \\
250,000 & 792.6 \(\pm 91.4\) & 0.893 & 0.924 & 0.966 & 0.135 & 51.6 &
860 \\
275,000 & 819.9 \(\pm 56.9\) & 0.920 & 0.934 & 0.971 & 0.120 & 52.9 &
854 \\
300,000 & 847.3 \(\pm 22.5\) & 0.947 & 0.944 & 0.975 & 0.106 & 54.3 &
848 \\
325,000 & 849.0 \(\pm 37.7\) & 0.947 & 0.951 & 0.979 & 0.108 & 54.8 &
835 \\
350,000 & 850.8 \(\pm 52.9\) & 0.948 & 0.957 & 0.983 & 0.110 & 55.2 &
821 \\
375,000 & 852.5 \(\pm 68.1\) & 0.948 & 0.964 & 0.986 & 0.111 & 55.7 &
808 \\
400,000 & 854.2 \(\pm 83.3\) & 0.948 & 0.970 & 0.990 & 0.113 & 56.1 &
794 \\
425,000 & 849.3 \(\pm 89.6\) & 0.943 & 0.967 & 0.992 & 0.115 & 56.3 &
787 \\
450,000 & 844.5 \(\pm 96.0\) & 0.938 & 0.964 & 0.994 & 0.117 & 56.4 &
779 \\
475,000 & 839.6 \(\pm 102.3\) & 0.933 & 0.961 & 0.995 & 0.119 & 56.6 &
772 \\
500,000 & 834.7 \(\pm 108.6\) & 0.928 & 0.958 & 0.997 & 0.121 & 56.7 &
764 \\
525,000 & 854.3 \(\pm 83.4\) & 0.944 & 0.964 & 0.999 & 0.114 & 58.8 &
756 \\
550,000 & 873.8 \(\pm 58.1\) & 0.960 & 0.970 & 1.000 & 0.106 & 60.9 &
748 \\
575,000 & 853.4 \(\pm 68.6\) & 0.947 & 0.969 & 0.991 & 0.125 & 59.3 &
762 \\
600,000 & 833.0 \(\pm 79.1\) & 0.933 & 0.968 & 0.982 & 0.144 & 57.7 &
775 \\
625,000 & 846.1 \(\pm 66.3\) & 0.944 & 0.964 & 0.987 & 0.140 & 58.5 &
761 \\
650,000 & 859.2 \(\pm 53.4\) & 0.954 & 0.960 & 0.992 & 0.136 & 59.3 &
748 \\
675,000 & 872.3 \(\pm 40.6\) & 0.965 & 0.956 & 0.997 & 0.132 & 60.1 &
734 \\
700,000 & 866.0 \(\pm\) 45.0 & 0.962 & 0.962 & 0.995 & 0.137 & 59.6 &
736 \\
725,000 & 860.2 \(\pm 28.1\) & 0.960 & 0.969 & 0.990 & 0.143 & 59.0 &
738 \\
750,000 & 870.4 \(\pm 48.0\) & 0.975 & 0.971 & 1.000 & 0.102 & 60.8 &
721 \\
775,000 & 874.6 \(\pm 29.0\) & 0.979 & 0.976 & 0.991 & 0.098 & 62.1 &
706 \\
800,000 & 876.5 \(\pm 29.1\) & 0.980 & 0.976 & 0.992 & 0.097 & 62.4 &
700 \\
825,000 & 878.0 \(\pm 29.2\) & 0.980 & 0.977 & 0.993 & 0.096 & 62.7 &
694 \\
850,000 & 872.9 \(\pm 30.4\) & 0.981 & 0.977 & 0.993 & 0.095 & 63.1 &
687 \\
875,000 & 879.2 \(\pm 29.4\) & 0.982 & 0.978 & 0.994 & 0.095 & 63.4 &
681 \\
900,000 & 874.1 \(\pm 31.2\) & 0.982 & 0.978 & 0.995 & 0.094 & 63.7 &
675 \\
925,000 & 877.6 \(\pm 28.7\) & 0.983 & 0.979 & 0.996 & 0.093 & 64.0 &
669 \\
950,000 & 871.8 \(\pm 30.1\) & 0.984 & 0.979 & 0.996 & 0.092 & 64.4 &
662 \\
975,000 & 879.4 \(\pm 29.9\) & 0.984 & 0.980 & 0.997 & 0.091 & 64.7 &
656 \\
1,000,000 & 875.7 \(\pm 30.0\) & 0.985 & 0.980 & 0.998 & 0.090 & 65.0 &
650 \\
\end{longtable}

\begin{figure}
\centering
\includegraphics[width=0.88\linewidth]{fig_return.png}
\caption{Return trajectory. Error bars show the reported evaluation standard deviation, not across-seed uncertainty.}
\end{figure}

\begin{figure}
\centering
\includegraphics[width=0.88\linewidth]{fig_coverage.png}
\caption{Mean and median coverage for the recorded runs. The no-CF series ends at 450,002 frames.}
\end{figure}

\begin{figure}
\centering
\includegraphics[width=0.88\linewidth]{fig_offroad.png}
\caption{Off-road fraction over training. The no-CF series ends at 450,002 frames.}
\end{figure}

\begin{figure}
\centering
\includegraphics[width=0.88\linewidth]{fig_speed.png}
\caption{Mean speed over training. The no-CF series ends at 450,002 frames.}
\end{figure}

\begin{longtable}[]{@{}
  >{\raggedleft\arraybackslash}p{(\columnwidth - 18\tabcolsep) * \real{0.1000}}
  >{\raggedleft\arraybackslash}p{(\columnwidth - 18\tabcolsep) * \real{0.1000}}
  >{\raggedleft\arraybackslash}p{(\columnwidth - 18\tabcolsep) * \real{0.1000}}
  >{\raggedleft\arraybackslash}p{(\columnwidth - 18\tabcolsep) * \real{0.1000}}
  >{\raggedleft\arraybackslash}p{(\columnwidth - 18\tabcolsep) * \real{0.1000}}
  >{\raggedleft\arraybackslash}p{(\columnwidth - 18\tabcolsep) * \real{0.1000}}
  >{\raggedleft\arraybackslash}p{(\columnwidth - 18\tabcolsep) * \real{0.1000}}
  >{\raggedleft\arraybackslash}p{(\columnwidth - 18\tabcolsep) * \real{0.1000}}
  >{\raggedleft\arraybackslash}p{(\columnwidth - 18\tabcolsep) * \real{0.1000}}
  >{\raggedleft\arraybackslash}p{(\columnwidth - 18\tabcolsep) * \real{0.1000}}@{}}
\caption{V-CARA without counterfactual evaluation}\label{tab:no-cf}\\
\toprule\noalign{}
\begin{minipage}[b]{\linewidth}\raggedleft
Frames
\end{minipage} & \begin{minipage}[b]{\linewidth}\raggedleft
Return \(\pm\) SD
\end{minipage} & \begin{minipage}[b]{\linewidth}\raggedleft
EMA
\end{minipage} & \begin{minipage}[b]{\linewidth}\raggedleft
Mean cov.
\end{minipage} & \begin{minipage}[b]{\linewidth}\raggedleft
Median
\end{minipage} & \begin{minipage}[b]{\linewidth}\raggedleft
P90
\end{minipage} & \begin{minipage}[b]{\linewidth}\raggedleft
Off-road
\end{minipage} & \begin{minipage}[b]{\linewidth}\raggedleft
OOF
\end{minipage} & \begin{minipage}[b]{\linewidth}\raggedleft
Speed
\end{minipage} & \begin{minipage}[b]{\linewidth}\raggedleft
T90
\end{minipage} \\
\midrule\noalign{}
\endhead
\bottomrule\noalign{}
\endlastfoot
25,000 & 110.9 \(\pm\) 74.1 & 110.9 & 0.22 & 0.22 & 0.29 & 0.65 & 0.20 &
32.5 & -1 \\
50,000 & 341.2 \(\pm\) 174.4 & 203.0 & 0.44 & 0.43 & 0.69 & 0.43 & 0.00
& 32.4 & -1 \\
75,001 & 491.8 \(\pm\) 243.5 & 318.5 & 0.59 & 0.58 & 0.94 & 0.19 & 0.00
& 35.4 & 866 \\
100,003 & 660.4 \(\pm\) 215.7 & 455.3 & 0.76 & 0.85 & 0.95 & 0.13 & 0.00
& 44.3 & 928 \\
125,000 & 838.4 \(\pm\) 58.1 & 608.5 & 0.93 & 0.93 & 0.97 & 0.05 & 0.00
& 52.2 & 862 \\
150,000 & 837.6 \(\pm\) 88.8 & 700.1 & 0.93 & 0.97 & 0.98 & 0.06 & 0.00
& 50.7 & 828 \\
175,002 & 743.5 \(\pm\) 181.2 & 717.5 & 0.84 & 0.90 & 0.95 & 0.19 & 0.00
& 48.3 & 818 \\
200,002 & 815.5 \(\pm\) 94.5 & 756.7 & 0.91 & 0.92 & 0.97 & 0.09 & 0.00
& 52.8 & 840 \\
225,002 & 830.5 \(\pm\) 48.2 & 786.2 & 0.93 & 0.94 & 0.97 & 0.15 & 0.00
& 56.7 & 791 \\
250,003 & 704.6 \(\pm\) 187.3 & 753.6 & 0.80 & 0.89 & 0.94 & 0.22 & 0.00
& 46.8 & 744 \\
275,003 & 831.5 \(\pm\) 29.2 & 784.8 & 0.93 & 0.93 & 0.95 & 0.12 & 0.00
& 53.7 & 758 \\
300,002 & 649.0 \(\pm\) 219.4 & 730.4 & 0.75 & 0.87 & 0.93 & 0.28 & 0.00
& 44.7 & 846 \\
325,000 & 697.9 \(\pm\) 236.1 & 717.4 & 0.79 & 0.89 & 0.98 & 0.24 & 0.00
& 47.1 & 751 \\
350,000 & 833.1 \(\pm\) 61.1 & 763.7 & 0.93 & 0.93 & 0.98 & 0.15 & 0.00
& 56.0 & 741 \\
375,003 & 844.6 \(\pm\) 54.6 & 796.1 & 0.94 & 0.94 & 0.98 & 0.13 & 0.00
& 54.9 & 792 \\
400,002 & 851.3 \(\pm\) 56.7 & 818.2 & 0.94 & 0.94 & 0.99 & 0.18 & 0.00
& 54.5 & 736 \\
425,000 & 786.5 \(\pm\) 194.3 & 805.5 & 0.88 & 0.95 & 0.97 & 0.18 & 0.00
& 54.7 & 750 \\
450,002 & 848.6 \(\pm\) 55.7 & 822.7 & 0.94 & 0.94 & 0.98 & 0.15 & 0.00
& 57.8 & 757 \\
\end{longtable}

\begin{longtable}[]{@{}
  >{\raggedleft\arraybackslash}p{(\columnwidth - 12\tabcolsep) * \real{0.1429}}
  >{\raggedleft\arraybackslash}p{(\columnwidth - 12\tabcolsep) * \real{0.1429}}
  >{\raggedleft\arraybackslash}p{(\columnwidth - 12\tabcolsep) * \real{0.1429}}
  >{\raggedleft\arraybackslash}p{(\columnwidth - 12\tabcolsep) * \real{0.1429}}
  >{\raggedleft\arraybackslash}p{(\columnwidth - 12\tabcolsep) * \real{0.1429}}
  >{\raggedleft\arraybackslash}p{(\columnwidth - 12\tabcolsep) * \real{0.1429}}
  >{\raggedleft\arraybackslash}p{(\columnwidth - 12\tabcolsep) * \real{0.1429}}@{}}
\caption{Common-window checkpoint comparison}\label{tab:common-window}\\
\toprule\noalign{}
\begin{minipage}[b]{\linewidth}\raggedleft
Approx. frames
\end{minipage} & \begin{minipage}[b]{\linewidth}\raggedleft
V-CARA return
\end{minipage} & \begin{minipage}[b]{\linewidth}\raggedleft
No-CF return
\end{minipage} & \begin{minipage}[b]{\linewidth}\raggedleft
V-CARA cov.
\end{minipage} & \begin{minipage}[b]{\linewidth}\raggedleft
No-CF cov.
\end{minipage} & \begin{minipage}[b]{\linewidth}\raggedleft
V-CARA off
\end{minipage} & \begin{minipage}[b]{\linewidth}\raggedleft
No-CF off
\end{minipage} \\
\midrule\noalign{}
\endhead
\bottomrule\noalign{}
\endlastfoot
100k & 427.0 \(\pm\) 134.7 & 660.4 \(\pm\) 215.7 & 0.527 & 0.76 & 0.110
& 0.13 \\
125k & 504.7 \(\pm\) 141.1 & 838.4 \(\pm\) 58.1 & 0.605 & 0.93 & 0.123 &
0.05 \\
150k & 582.4 \(\pm\) 147.5 & 837.6 \(\pm\) 88.8 & 0.683 & 0.93 & 0.137 &
0.06 \\
175k & 660.1 \(\pm\) 153.8 & 743.5 \(\pm\) 181.2 & 0.760 & 0.84 & 0.150
& 0.19 \\
200k & 737.8 \(\pm\) 160.2 & 815.5 \(\pm\) 94.5 & 0.838 & 0.91 & 0.163 &
0.09 \\
225k & 765.2 \(\pm\) 125.8 & 830.5 \(\pm\) 48.2 & 0.865 & 0.93 & 0.149 &
0.15 \\
250k & 792.6 \(\pm\) 91.4 & 704.6 \(\pm\) 187.3 & 0.893 & 0.80 & 0.135 &
0.22 \\
275k & 819.9 \(\pm\) 56.9 & 831.5 \(\pm\) 29.2 & 0.920 & 0.93 & 0.120 &
0.12 \\
300k & 847.3 \(\pm\) 22.5 & 649.0 \(\pm\) 219.4 & 0.947 & 0.75 & 0.106 &
0.28 \\
325k & 849.0 \(\pm\) 37.7 & 697.9 \(\pm\) 236.1 & 0.947 & 0.79 & 0.108 &
0.24 \\
350k & 850.8 \(\pm\) 52.9 & 833.1 \(\pm\) 61.1 & 0.948 & 0.93 & 0.110 &
0.15 \\
375k & 852.5 \(\pm\) 68.1 & 844.6 \(\pm\) 54.6 & 0.948 & 0.94 & 0.111 &
0.13 \\
400k & 854.2 \(\pm\) 83.3 & 851.3 \(\pm\) 56.7 & 0.948 & 0.94 & 0.113 &
0.18 \\
425k & 849.3 \(\pm\) 89.6 & 786.5 \(\pm\) 194.3 & 0.943 & 0.88 & 0.115 &
0.18 \\
450k & 844.5 \(\pm\) 96.0 & 848.6 \(\pm\) 55.7 & 0.938 & 0.94 & 0.117 &
0.15 \\
\end{longtable}

\begin{figure}
\centering
\includegraphics[width=0.88\linewidth]{fig_common_window_delta.png}
\caption{Descriptive differences in mean coverage and off-road fraction at nearest logged checkpoints within 25,000 frames. Values are V-CARA minus no-CF; this is not a paired statistical test.}
\label{fig:common-window-delta}
\end{figure}

\section{Discussion}\label{discussion}

The data show why return, coverage, speed, and off-road behaviour should
be considered together. In the full V-CARA trajectory, mean speed generally rises with coverage.
The 700k checkpoint shows a modest decline relative to 675k rather than a
large collapse; the larger regression is concentrated around 425k--500k. In the no-CF trajectory, speed increases early but falls at
the same checkpoints where coverage and return regress. These correlated
changes suggest instability or policy degradation rather than a simple
monotonic learning process.

The most visible difference between the two runs occurs around
250k--325k frames. Full V-CARA has lower off-road fraction and higher
mean coverage at 250k, 300k, and 325k. However, the data do not isolate
which internal component causes this pattern. Counterfactual evaluation
is the intended difference, but stochastic training, model learning,
replay contents, and checkpoint noise can also contribute. Multiple
seeds and logged counterfactual acceptance statistics are needed for a
causal claim.

The V2-to-current design history emphasizes implementation semantics. An
uncertainty threshold on the wrong scale can disable counterfactual
planning; replay buffers conditioned on the wrong terminal flag can
discard useful trajectories; an auxiliary head with weak labels can
contribute little or destabilize actor updates; and a brake dead zone
can materially alter the control distribution. These are not minor
engineering details: they define the algorithm that was actually
trained.

Across the late full V-CARA checkpoints from 750k to 1M frames, mean
return ranges from 870.4 to 879.4, while mean coverage increases from 0.975
to 0.985 and off-road fraction decreases from 0.102 to 0.090. The return
remains within a relatively narrow range as coverage and off-road behaviour
continue to change; these observations alone do not establish saturation
of the return objective.

\section{Limitations and Threats to
Validity}\label{limitations-and-threats-to-validity}

First, each condition is represented by a single seed-42 training
trajectory. The standard deviations in return describe variation within
the reported evaluation episodes, not variation across independent
training seeds. They should not be interpreted as confidence intervals
for algorithm-level performance.

Second, the evaluation series have different lengths. Full V-CARA
extends to one million frames, while no-CF ends at 450,002. Thus the
direct comparison is restricted to the common window. No extrapolation
of no-CF performance beyond its last observed checkpoint is made.

Third, the coverage and off-road metrics depend on the environment's
implementation. A reproducible paper should define exactly how coverage
is normalized, how off-road frames are counted, how OOF differs from
off-road, and how T90 is computed. The current manuscript uses the names
and values in the recorded evaluation logs.

Fourth, the reported results are checkpoint-level summaries from one training seed per condition. Raw per-episode evaluation records and a run manifest should accompany a public reproducibility release so that each summary can be regenerated from its source data. Fifth, no held-out environment, track variation, or cross-domain evaluation is included, so generalization beyond the evaluated task is not established.

Finally, the paper describes the implemented architecture and observed
outcomes but does not claim statistical significance or universal
superiority. A confirmatory experiment should use multiple training
seeds, identical frame budgets, matched evaluation seeds, and
predeclared primary outcomes.

\section{Reproducibility and Future
Work}\label{reproducibility-and-future-work}

A reproducibility release should include the environment version, code
commit, complete configuration, random seeds, model checkpoints, raw
per-episode metrics, and scripts used to generate every figure and
table. The configuration should specify image size, frame stack, action
repeat, time limit, stall rules, reward scaling, actor/critic learning
rates, replay capacity and mixture, counterfactual candidate set,
uncertainty normalization, risk weights, and all auxiliary-loss gates.

The next experiment should run full V-CARA and no-CF for equal budgets
using at least five independent training seeds per condition, with a
fixed evaluation suite. Report means and standard deviations across
training seeds separately from episode-level uncertainty. A paired
analysis should use the same evaluation seed IDs and report per-seed
differences for return, mean coverage, P90 coverage, off-road fraction,
speed, and T90.

Ablations should remove one mechanism at a time: counterfactual action
evaluation, uncertainty gating, risk scoring, future-coverage
prediction, and coverage-aware replay. Runtime should be measured
alongside environment frames and performance, since the latent ensemble
and candidate scoring add compute. The decisive question is not only
whether the final score improves, but whether the improvement is
repeatable, interpretable, and worth the additional inference and
training cost.

\section{Conclusion}\label{conclusion}

This paper presents V-CARA, a visual reinforcement-learning method that
augments SAC with short-horizon latent dynamics, uncertainty-gated local
counterfactual action evaluation, risk and progress prediction,
future-coverage supervision, and coverage-aware replay. On the full
V-CARA seed-42 trajectory, the one-million-frame checkpoint reaches
mean coverage 0.985, median 0.980, P90 0.998, return 875.7 \(\pm\) 30.0,
off-road fraction 0.090, mean speed 65.0, and T90=650.

The comparison with the no-CF run is not uniformly one-sided. No-CF
progresses faster in the early checkpoints, while full V-CARA shows
markedly higher coverage and lower off-road fraction in the 250k--325k
interval. The evidence supports a nuanced finding: counterfactual
evaluation may help stabilize a difficult middle phase in this
trajectory, but its causal contribution and across-seed reliability
remain to be established. This is the basis for a testable research
program rather than a claim of universal superiority.

\section{References}\label{references}

\begin{enumerate}
\def\labelenumi{\arabic{enumi}.}
\tightlist
\item
  Haarnoja, T., Zhou, A., Abbeel, P., and Levine, S. (2018). Soft
  Actor-Critic: Off-Policy Maximum Entropy Deep Reinforcement Learning
  with a Stochastic Actor. \emph{Proceedings of the 35th International
  Conference on Machine Learning}, 80, 1861--1870.
\item
  Janner, M., Fu, J., Zhang, M., and Levine, S. (2019). When to Trust
  Your Model: Model-Based Policy Optimization. \emph{Advances in Neural
  Information Processing Systems}, 32.
\item
  Schaul, T., Quan, J., Antonoglou, I., and Silver, D. (2016).
  Prioritized Experience Replay. \emph{International Conference on
  Learning Representations}.
\item
  Henderson, P., Islam, R., Bachman, P., Pineau, J., Precup, D., and
  Meger, D. (2018). Deep Reinforcement Learning that Matters.
  \emph{Proceedings of the AAAI Conference on Artificial Intelligence},
  32(1).
\item
  Sutton, R. S., and Barto, A. G. (2018). \emph{Reinforcement Learning:
  An Introduction}, 2nd edition. MIT Press.
\end{enumerate}

\section{Data and Reproducibility
Statement}\label{data-and-reproducibility-statement}

The reported tables summarize the V-CARA seed-42 evaluation series through
1,000,000 frames and the V-CARA-no-CF seed-42 trajectory through 450,002
frames. No missing checkpoint values are interpolated. Reproduction of the
results requires raw per-episode evaluations, checkpoint identifiers, the
source-code commit, environment version, complete configuration, and scripts
for generating all tables and figures.


\end{document}
